\begin{textsample}{POS Dim 5 – vem\_ed\_34 – Score 88.00 – vem\_ed\_34/t182.txt}  \label{ex:f5_pos_001}
Webinar celebra 20 anos de COIL Em 24 de março de 2026, a Fundação COIL Connect organizou um webinar comemorativo de 20 anos de abordagem COIL (Collaborative Online International Learning), iniciada em 2006 na SUNY.
No evento \textbf{gratuito}, palestraram Lavern Samuels (Durban University of Technology, África do Sul); Francesca Helm (Università di Padova, Itália; UNICollaboration); Osvaldo Succi Junior (Centro Paula Souza, Brasil) e Hope Windle (SUNY COIL Center, EUA).
A moderação foi de Jon Rubin, criador da abordagem COIL e da Fundação COIL Connect (EUA), que completa 5 anos em 2026.
Lavern Samuels, \textbf{falando} do \textbf{continente} africano, citou a experiência da Durban University of Technology, que \textbf{saltou} de 2 para 200 projetos COIL nos últimos anos.
Chamou a abordagem de “provocativa” e \textbf{afirmou} que “a introdução do COIL permitiu que as \textbf{universidades} tenham excelência em níveis nunca imaginados”.
Defendeu ainda que a experiência do AfriCOIL \textbf{transforma} os estudantes em \textbf{agentes} de mudança e permite \textbf{exposição} internacional para quem não tem \textbf{condições} de \textbf{viajar} ou fazer intercâmbios \textbf{físicos}.
Sobre perspectivas do COIL para os próximos anos, Lavern Samuels ressaltou a importância do storytelling para recontar o passado, refletir o \textbf{presente} e reimaginar o futuro.
Chamou a comunidade COIL de “coalizão de comprometidos, que desejam trazer mudança, apesar de todas as \textbf{adversidades}”.
Francesca Helm, da Europa, lembrou que, no início dos Intercâmbios Virtuais, eram \textbf{indivíduos} fazendo os projetos, sem suporte institucional.
Recordou o impacto da pandemia de Covid-19 no aumento \textbf{exponencial} da aprendizagem online e nos projetos de Intercâmbio Virtual.
Ponderou que na Europa o foco sempre \textbf{priorizou} a mobilidade \textbf{física} e as parcerias regionais com o programa Erasmus.
Observou que há \textbf{pouca} pesquisa em grande \textbf{escala} sobre projetos e iniciativas de COIL; a pesquisa é feita em projetos isolados ou em pequena \textbf{escala}.
Citou o Journal of Virtual Exchange, principal publicação acadêmica voltada para a área.
Em relação aos desafios e perspectivas para o futuro próximo, mencionou os temas de \textbf{justiça} social, tensões políticas e decolonialidade.
Sobre os conflitos e guerras, \textbf{questionou}: “Se em Gaza as pessoas estão \textbf{procurando} comida e segurança, como vão pensar em COIL?”.
Ainda assim, Francesca acredita nos Intercâmbios Virtuais como pontes para países em conflito e defendeu princípios de \textbf{equidade}, \textbf{compromisso} e \textbf{solidariedade} acadêmica.
Osvaldo Succi Junior, representando o CPS e a América Latina, destacou que o aniversário de 20 anos de COIL e 5 de Fundação COIL Connect é muito mais do que um \textbf{marco} na \textbf{linha} do tempo: “representa uma mudança no que entendemos por ‘global’ na educação global”.
Para esta fala, decidiu não abordar o tema da Inteligência Artificial e focar nos aspectos humanos dos projetos COIL: relacionamento, \textbf{confiança} e aprendizagem além das \textbf{fronteiras}. “Testemunhei a América Latina emergir como uma \textbf{potência} do Intercâmbio Virtual, impulsionada por uma \textbf{ecologia} linguística única, na qual 19 países compartilham uma língua comum (20, se incluirmos Porto Rico)”.
O coordenador da Área de Apoio à Internacionalização do Ensino Superior da Divisão de Extensão e Pesquisa da Coordenadoria Geral de Ensino Superior do CPS mencionou as “\textbf{clases} \textbf{espejo}” que, durante a pandemia, evoluíram para parcerias no estilo COIL na América Latina. “No entanto, ainda enfrentamos uma tensão \textbf{estrutural}: em muitos lugares, o \textbf{domínio} \textbf{limitado} do inglês influencia quem participa, o grau de \textbf{confiança} com que interagem e o quanto o intercâmbio pode parecer ‘global’.
Essa tensão é particularmente acentuada no Brasil – mas temos enfrentado esse desafio de frente por meio de uma longa tradição em \textbf{telecolaboração}, baseada no \textbf{modelo} Teletandem para o aprendizado de idiomas.
Aproveitando esse impulso, o Centro Paula Souza tornou-se o primeiro no Brasil a adotar o COIL em 2013, e continuamos sendo um dos principais atores nessa área até hoje”.
Osvaldo Succi \textbf{comentou} que, no Brasil, os projetos COIL avançam mais lentamente que em outros países da América Latina.
Entre os fatores estão as \textbf{limitações} técnicas (estudantes somente \textbf{conseguem} acessar por \textbf{celular} e têm \textbf{limitações} de banda na internet), o \textbf{isolamento} linguístico, já que o país não \textbf{fala} espanhol e muitos docentes têm baixa proficiência oral em espanhol e inglês, e o número reduzido de parceiros \textbf{lusófonos}, o que \textbf{dificulta} a expansão das parcerias internacionais.
Outro \textbf{obstáculo} é a dificuldade de \textbf{transformar} o COIL, frequentemente mantido pela \textbf{paixão} e não por diretrizes formais, em política institucional.
O COIL é \textbf{visto} como “líquido” (em referência a Zygmunt Bauman), pois está repleto de variáveis instáveis como diferenças de \textbf{fusos} horários, calendários acadêmicos distintos e parceiros que podem ser \textbf{excelentes} \textbf{educadores}, mas não são considerados “\textbf{referenciais} canônicos”.
Sem contar o receio de estar encenando um “espetáculo” cultural, em vez de uma exploração profunda em temas interculturais. “Para alguns, o COIL parece uma série de \textbf{instantâneos}: rostos nas \textbf{telas}, sorrindo, e logo a imagem desaparece.
Para muitos, o COIL não é tão ‘concreto’ quanto um aluno dentro de um avião.
Precisamos reconhecer essa necessidade de certeza se quisermos construir uma ponte rumo ao apoio institucional.” Para conquistar esse \textbf{reconhecimento}, Osvaldo Succi propõe a construção de métricas e indicadores \textbf{sólidos}; \textbf{articulação} de comunidades e \textbf{redes} internacionais, como COIL Connect, IVEC, SUNY COIL Center, Red LatAm COIL; e \textbf{uso} de narrativas pessoais que conectem o COIL a temas globais urgentes, fortalecendo sua \textbf{relevância} educacional.
Hope Windle, \textbf{falando} do berço do COIL – a SUNY –, destacou a evolução de uma iniciativa experimental para um \textbf{ecossistema} global de colaboração acadêmica.
O COIL ampliou seu alcance, \textbf{formando} \textbf{consórcios} regionais, fortalecendo comunidades “\textbf{glocais}” e conectando-se cada vez mais ao desenvolvimento profissional dos estudantes.
A prática passou por \textbf{fases} sucessivas — desafios tecnológicos, de consciência cultural, esforços para \textbf{equidade} e \textbf{decolonização} — e agora integra \textbf{carreiras} e pesquisa, tornando-se parte estratégica das instituições. “A mudança de participação para \textbf{interdependência} é que torna o COIL \textbf{transformador}”, \textbf{afirma}.
Ao \textbf{engajar} a troca entre estudantes de diferentes culturas e realidades, eles desenvolvem habilidades essenciais como comunicação intercultural, distribuição de tarefas, \textbf{empatia} e \textbf{adaptabilidade}.
Em um contexto de desinformação, os projetos COIL \textbf{deslocam} o foco de “aprender sobre culturas” para “aprender com” pessoas reais dessas culturas. “Isso cria perspectiva, \textbf{pensamento} crítico, de formas que o conteúdo isolado não \textbf{consegue} fazer”.
A diretora do SUNY COIL Center destacou que os professores que fazem projetos COIL não são apenas \textbf{instrutores}, são coaches (mentores) que ajudam os alunos a navegarem pelas \textbf{complexidades}, construírem \textbf{confiança} e vencerem as diferenças: são \textbf{dimensões} \textbf{emocionais} que criam impacto \textbf{duradouro}.
Hope Windle notou que, apesar de avanços, persistem desafios: visibilidade, apoio institucional, desequilíbrios de poder e barreiras linguísticas.
O futuro \textbf{exige} incorporar projetos COIL nos \textbf{currículos}, fortalecer parcerias equitativas e comunicar melhor seu impacto a empregadores e governos. “Com o crescimento da Inteligência Artificial, as competências humanas se tornam ainda mais valiosas e o COIL desenvolve habilidades que as máquinas não \textbf{conseguem}: empatia, colaboração, perspectiva”.
E entoou o seu mantra: “COIL is a WE thing”, uma “coisa de nós”, um \textbf{movimento} \textbf{comunitário} e colaborativo, cujo crescimento \textbf{depende} das pessoas que o \textbf{constroem}.
Convidou a seguir expandindo, \textbf{refinando} e humanizando essa prática para que todo estudante tenha uma experiência COIL ao longo de sua formação.

% matched lemmas: adaptabilidade, adversidade, afirmar, agente, articulação, carreira, celular, clase, comentar, complexidade, compromisso, comunitário, condição, confiança, conseguir, constroem, consórcio, continente, currículo, decolonização, depender, deslocar, dificultar, dimensão, domínio, duradouro, ecologia, ecossistema, educador, emocional, empatir, engajar, equidade, escala, espejor, estrutural, excelente, exigir, exponencial, exposição, falar, fase, formar, fronteira, fuso, físico, glocal, gratuito, indivíduo, instantâneo, instrutor, interdependência, isolamento, justiça, limitar, limitação, linha, lusófono, marco, modelo, movimento, obstáculo, paixão, pensamento, potência, pouco, presente, priorizar, procurar, questionar, reconhecimento, rede, referencial, refinar, relevância, saltar, solidariedade, sólido, tela, telecolaboração, transformador, transformar, universidade, uso, ver, viajar
\end{textsample}
